\id{ҒТАМР 52.01.93}{}

\begin{header}
\swa{}{ӨНДІРІСТІК ЖАРАҚАТТАНУ ҚАЗАҚСТАН РЕСПУБЛИКАСЫНЫҢ ФОСФОР ӨНЕРКӘСІБІ КӘСІПОРЫНДАРЫНДАҒЫ ЕҢБЕКТІ ҚОРҒАУ ЖӘНЕ КӘСІПТІК ТӘУЕКЕЛДЕРДІ БАСҚАРУ ЖАЙ КҮЙІНІҢ КӨРСЕТКІШІ РЕТІНДЕ}

Л.М.~Актаева,
А.Б.~Бекмагамбетов,
А.М.~Рахметова,
Г.А.~Насырова,
Э.А.~Кұлмағамбетова,
Г.К.~Алшынбекова,
А.Т.~Шілменова\envelope,
Н.Т.~Сағындықова
\end{header}

\begin{affil}
Қазақстан Республикасы Еңбек және халықты әлеуметтік қорғау министрлігінің Еңбекті қорғау жөніндегі республикалық ғылыми-зерттеу институты ШМҚ РМК, Астана, Қазақстан

\corrauthor{Корреспондент-автор: aigul.shilmenova@gmail.com}
\end{affil}

Мақалада еңбек жағдайларын бағалау мен кәсіптік тәуекелдерді басқару
тұрғысынан 2024 жылы фосфор өнеркісібі кәсіпорындарындағы жарақаттанудың
жай-күйі қарастырылған. Өндірістік жарақаттану санитарлық-гигиеналық
еңбек жағдайларын, кәсіптік қызметтің ерекшелігін және жұмысшының еңбек
процесінің субъектісі ретіндегі жағдайын көрсететін кәсіптік тәуекелдің
негізгі өлшемдерінің бірі ретінде қараастырылады. Зерттеудің мақсаты -
жазатайым оқиғаларды тіркеудің ерекшеліктерін анықтау және қауіпсіздік
көрсеткіштерінің сенімділігіне әсер ететін факторларды анықтау.
Әдістемелік негіз-жазатайым оқиғалар туралы ресми есептілік пен құжаттық
материалдарды талдау. Тіркелген 65 оқиғаның тек аз ғана бөлігі жұмыс
міндеттерін орындаумен байланыста екені анықталды, бұл тіркеудің
жеткіліксіздігін және кәсіби тәуекелдердің нақты деңгейінің
бұрмаланғанын көрсетеді. Қызметкерлердің жасын, еңбек өтілін және кәсіби
тәжірибесін ескере отырып, бес өндірістік бөлімшедегі оқиғалардың
салыстырмалы талдауы жүргізілді. Нәтижелер қолданыстағы бақылау
жүйесінде еңбек қауіпсіздігін объективті бағалау мүмкіндігінің шектеулі
екенін көрсетеді. Еңбек қауіпсіздігін қамтамасыз етудің реттеуші және
ұйымдастырушылық тетіктерін жетілдіру және тәуекелге негізделген тәсілді
енгізу қажеттілігі туралы қорытынды жасалды.

{\bfseries Түйін сөздер:} еңбек жағдайлары, еңбек қауіпсіздігі және еңбекті
қорғау, өндірістік жарақаттану, тұрмыстық жарақаттану, жазатайым
оқиғалар, кәсіби тәуекелділік.

\begin{header}
ПРОИЗВОДСТВЕННЫЙ ТРАВМАТИЗМ КАК ПОКАЗАТЕЛЬ СОСТОЯНИЯ ОХРАНЫ ТРУДА И УПРАВЛЕНИЯ ПРОФЕССИОНАЛЬНЫМИ РИСКАМИ НА ПРЕДПРИЯТИЯХ ФОСФОРНОЙ ПРОМЫШЛЕННОСТИ РЕСПУБЛИКИ КАЗАХСТАН

Л.М.~Актаева,
А.Б.~Бекмагамбетов,
А.М.~Рахметова,
Г.А.~Насырова,
Э.А.~Кульмагамбетова,
Г.К.~Алшынбекова,
А.Т.~Шилменова\envelope,
Н.Т.~Сагиндикова
\end{header}

\begin{affil}
РГП на ПХВ Республиканский научно-исследовательский институт по охране труда Министерства труда и социальной защите населения Республики Казахстан, Астана, Казахстан,

e-mail: aigul.shilmenova@gmail.com
\end{affil}

В статье рассмотрено состояние травматизма на предприятиях фосфорной
промышленности в 2024 г. С позиции оценки условий труда и управления
профессиональными рисками. Производственный травматизм рассматривается
как один из ключевых критериев профессионального риска, отражающий
санитарно-гигиенические условия труда, специфику профессиональной
деятельности и состояние работника как субъекта трудового процесса.
Целью исследования заключается в выявлении особенностей учета несчастных
случаев и определения факторов, влияющих на достоверность показателей
безопасности. Методологической основой послужили анализ официальной
отчетности и документальных материалов по несчастным случаям.
Установлено, что при общем количестве 65 зарегистрированных происшествий
лишь незначительная часть связана с выполнением трудовых обязанностей,
что свидетельствует о недостаточной полноте регистрации и искажении
реального уровня профессиональных рисков. Проведен сравнительный анализ
происшествий в пяти производственных подразделениях с учетом возраста,
стажа и профессиональной принадлежности работников. Полученные
результаты указывают на ограниченные возможности объективной оценки
состояния охраны труда в рамках действующей системы контроля. Сделан
вывод о необходимости совершенствования нормативно правовых и
организационных механизмов обеспечения безопасности труда и внедрения
риск-ориентированного подхода.

{\bfseries Ключевые слова:} условия труда, безопасность и охрана труда,
производственный травматизм, бытовой травматизм, несчастные случаи,
профессиональный риск.

\begin{header}
OCCUPATIONAL INJURIES AS AN INDICATOR OF THE STATE OF OCCUPATIONAL SAFETY AND PROFESSIONAL RISK MANAGEMENT AT ENTERPRISES OF THE PHOSPHORUS INDUSTRY OF THE REPUBLIC OF KAZAKHSTAN

L.~MAktayeva,
A.B.~Bekmagambetov,
А.M.~Rakhmetova,
G.A.~Nasyrova,
E.A.~Kulmagambetova,
G.A.~Alshynbekova,
A.T~Shilmenova\envelope,
N.T.~Sagindikova
\end{header}

\begin{affil}
RSE at the National Research Institute for Occupational Safety of the Ministry of Labor and Social Protection of the Population of the Republic of Kazakhstan, Astana, Kazakhstan,

e-mail: aigul.shilmenova@gmail.com
\end{affil}

This article examines the injury rate at phosphorus industry enterprises
in 2024 from the perspective of assessing working conditions and
managing occupational risks. Occupational injuries are considered a key
criterion for occupational risk, reflecting sanitary and hygienic
working conditions, the specific nature of professional activity and the
worker's status as a subject of the labor process7 the aim of the study
is to identify the characteristics of accident recording and determine
the factors affecting the reliability of safety indicators. The
methodological basis was an analysis of official reports and documentary
materials on accidents. It was found that, out of a total of 65
registered incidents, only a small proportion were related to the
performance of work duties, indicating insufficient registration and a
distortion of the actual level of occupational risks. A comparative
analysis of incidents in five production units was conducted,
considering the age, length of service and professional affiliation of
workers. The results obtained indicate the limited possibilities for an
objective assessment of the state of occupational safety within the
framework of the current control system. The study concludes that it is
necessary to improve regulatory and organiza\-tion mechanisms for ensuring
occupational safety and implement a risk-based approach.

{\bfseries Keywords:} working conditions, occupational safety and health,
industrial injuries, domestic injuries, accidents, occupational risks.

\begin{multicols}{2}
{\bfseries Кіріспе.} Өндірістік жарақаттар еңбек жағдайларының жай-күйін
және кәсіпорындардағы кәсіптік тәуекелдерді басқарудың тиімділігін
көрсететін негізгі еңбек қауіпсіздігі мәселесі. Өндірістік жарақат -
жұмысшының еңбек міндеттерін орындау кезінде денсаулығына келтірілген
және уақытша немесе тұрақты мүгедектікке әкеп соққан зардап ретінде
анықталады {[}1{]}. Өндірістік жарақаттардың жиілігі өндірістік
үдерістердің қауіпсіздігінің және еңбекті қорғау жүйесінің жұмыс істеу
сапасының ажырамас көрсеткіші болып табылады.

Өндірістік жарақаттануды зерттеу еңбек қауіпсіздігі мен еңбекті
қорғаудың, қоғамдық денсаулық сақтаудың және әлеуметтік-экономикалық
дамудың өзекті мәселесі болып табылады. Өйткені өндірістегі жазатайым
оқиғалар айтарлықтай адами және экономикалық шығындармен қатар жүреді.
Бұл мәселе әсіресе өндірістік қауіптілігі жоғары салаларда, зиянды және
қауіпті өндіріс жағдайында анық көрінеді {[}2-4{]}.

Жазатайым оқиғалар мен кәсіптік аурулардың салдарынан Халықаралық еңбек
ұйымының (бұдан әрі - ХЕҰ) деректері бойынша әлемде 3 миллионға жуық
адам қаза табады, ал өндірістегі жазатайым оқиғалардың жалпы саны жылына
374 млн-нан асады {[}2{]}. Сонымен қатар, әлемнің көптеген елдерінің
ғалымдарының зерттеулері өндірістік орта тудыратын зиянды және қауіпті
факторлар денсаулықтың ауытқуы себеп-салдарының 30\%-на дейін әсер
ететінін көрсетеді {[}3,4{]}. Бұл деректер мәселенің жаһандық сипатын,
әсіресе өндірістік қауіптілігі жоғары салаларда кәсіби тәуекелдерді
басқару тәсілдерін жетілдіру қажеттілігін құптайды {[}5{]}. ДДҰ
мәліметтері бойынша әлемнің көптеген елдерінде «жұмыспен байланысты
денсаулық мәселері себебінен» болған экономикалық шығындар ЖҰӨ-нің 4-6\%
құрайды {[}6{]}.

Өндірістік жарақаттану деңгейі еңбек процесін ұйымдастырудың қазіргі
жағдайы мен еңбек сипатына (қауіпті, ауыр, жүйке-күйзелістік, ауысымды
жұмыс кестесі) байланысты. Бұл көрсеткіш жаңа қауіпсіз технологиялар мен
жабдықтарды енгізуге, өндіріс мәдениетін арттыруға және тағы да басқа
өзгерістерге жылдам әрекет етеді, басқаша айтқанда, өндірістік жарақат
еңбек жағдайлары мен еңбекті қорғаудың жай-күйінің сезімтал индикаторы
ретінде еңбекті қорғауды басқарудың жаңа тәсілі негізделген кәсіби
тәуекелді бағалаудың маңызды өлшемшарты болуы тиіс.

Халықаралық тәжірибеде өндірістік жарақаттану деңгейін бағалау және
салыстыру үшін салыстырмалы көрсеткіштер, соның ішінде 100~000
жұмыскерге шаққанда зардап шеккендер мен қаза тапқандар жиілігінің
коэффициенті қолданылады.

Қазақстанда 2022 жылдың қорытындысы бойынша жазатайым оқиғалардан зардап
шеккендердің жиілік коэффициенті 100~000 қызметкерге 44 адамды, қаза
тапқандардың жиілік коэффициенті-3,7 (2020 жылы-тиісінше 39 және 3,9)
құрады. Салыстыру үшін, дамыған елдерде зардап шеккендердің жиілік
коэффициенті едәуір жоғары: АҚШ- 900 және 5,3, Италия-1413 және 3,4,
Франция-2931 және 2,5, бұл ретте өлімге әкелетін жарақаттану
көрсеткіштері өз елімізбен салыстырғанда бірдей немесе төмен көрсеткішке
ие {[}7,8{]}.

Өндірістегі жазатайым оқиғалардың Еуропалық статистикасында {[}9{]},
2018 жылы Еуро Одақ елдерінде шамамен 3,3 миллион өлімге алып келмейтін
жарақаттар орын алды, бұл кем дегенде төрт күнтізбелік күн жұмыссыз
қалуға алып келді. АҚШ-та Ұлттық қауіпсіздік кеңесінің мәліметтері
бойынша, 2019 жылы 4,65 миллион өндірістік жарақат тіркелді, нәтижесінде
105 миллион өндірістік күн жоғалды және өндірістік жарақаттарды емдеудің
жалпы құны 171 миллиард долларға бағаланған. Еуропалық Одақ елдерінде,
шетелдік авторлардың мәліметтері бойынша, жыл сайын жұмыс орындарында
өндіріске байланысты себептер бойынша 350 мыңға жуық адам қаза табады
{[}10{]}. Ресейде 2022 жылы өндірісте ауыр зардаптармен 4639 жазатайым
оқиға болды, оның ішінде 991 адам қайтыс болды {[}11,12{]}.

Қазақстан Республикасы Ұлттық статистика бюросының деректеріне сәйкес,
экономика саласында бір және одан да көп жұмыс күніне шаққанда еңбек
қабілеттілігінен айырылған зардап шегушілердің ең көп саны, оның ішінде
өлім-жітім жағдайлары, өңдеу өнеркәсібінде 800 жағдай тіркелді, бұл
жазатайым оқиғалардың жалпы санының 32,4\% құрайды, оның ішінде 50
топтық жазатайым оқиға орын алды. Бұл салада 36 өлім тіркелді, сондай-ақ
ауыр жарақат басым. Зардап шеккендердің ең көп саны Қарағанды, Ұлытау
және Шығыс Қазақстан облыстарында байқалады, бұл өнеркәсіптік
кәсіпорындардың тығыз шоғырлануына байланысты {[}13{]}.

Өңдеу секторына жататын фосфор өнеркәсібінің кәсіпорындары күрделі,
технологиялық үрдістермен, химиялық қауіпті фосфор қосылыстарының
әсерімен, өрт пен жарылыс қаупінің жоғарылыуымен, сондай-ақ өндірістік
үдерістерді механикаландыру мен автоматтандырудың жеткіліксіз деңгейімен
сипатталады. Бұл жағдайда өндірістік жарақат еңбекті қорғау жағдайы мен
кәсіби тәуекелдерді басқару тиімділігінің маңызды көрсеткіші болып
табылады.

Осы зерттеуде фосфор өнеркәсібіндегі өндірістік жарақаттану, жазатайым
оқиғалардың жасына, жұмыс өтіліне, кәсібіне және жарақаттарды
оқшаулануына қарай бөлінуін ескере отырып талданады, бұл жұмысшылардың
ең осал топтарын және өндірістің маңызды аймақтарын анықтауға мүмкіндік
береді.

Зерттеудің мақсаты-еңбекті қорғау және кәсіптік тәуекелдерді басқару
жай-күйінің индикаторы ретінде фосфор өнеркәсібі кәсіпорындарындағы
өндірістік жарақаттануды талдау, жарақаттануға әсер ететін факторларды
анықтау және жазатайым оқиғалардың туындау қаупін бағалау болып
табылады.

{\bfseries Материалдар мен әдістер.} Зерттеу сипаттамалық құжаттық сипатта,
2024 жылдың қаңтар-желтоқсан айлары аралығында жүргізілді. Зерттеу
объектісі өңдеу саласындағы пилоттық кәсіпорынның бес өндірістік
цехының, оның ішінде ұсақтау-кептіру, сары фосфор алуда фосфорит
агломератын өндіру, сары фосфор өндіру, фосфорды қайта түзу,
бейтараптандыру мен жағу, азотты-оттегі өндірісінің 610 қызметкері
болды.

Қатысушылар құжаттық іріктеу әдісімен таңдалды. Зардап шеккендердің
санын анықтау үшін еңбекке уақытша жарамсыздық парақтары пайдаланылды,
ал кәсіби тәуекелді бағалау барлық цехтар арасында жүргізілді.
Сипаттамалық талдау үшін 2024 жылғы Н-1 жазатайым оқиға ("Еңбек
қызметіне байланысты жазатайым оқиғаларды тергеп-тексеру материалдарын
ресімдеу бойынша нысандарды бекіту туралы" Қазақстан Республикасы
Денсаулық сақтау және әлеуметтік даму министрінің 2015 жылғы 28
желтоқсандағы № 1055 бұйрығы) және 7-ТПЗ нысандары бойынша «Еңбек
қызметімен байланысты жарақаттану және кәсіптік аурулар туралы есеп»
ресми статистикалық деректер қолданылды.

{\bfseries Нәтижелер және талқылаулар. З}ерттеу пилоттық кәсіпорында
жүргізілді. Жұмысшылардың жалпы саны 45 мамандық бойынша 1162 адамды
құрады.

Зерттеулер пилоттық кәсіпорында толық 2024 жылдағы еңбекке жарамсыздық
парақтары бойынша жүргізілді. Тұрмыстық жарақаттанудың 59 жағдайын қоса
алғанда, 610 еңбекке жарамсыздық парақтары зерттелді.

\emph{Ұсақтау-кептіру цехы} шикізатты дайындауға (шикізатты қабылдау,
қоймалау, ұсақтау, кептіру және сұрыптау), содан кейін оны пеш цехы мен
агломерация цехының дестелеу бөлімшесіне жіберу үшін бағытталған.

Цехта негізгі технологиялық құрал-жабдықтар қолданылады, оларға
денсаулыққа қауіпті шикізат көздерімен қатар жүретін роторлы
стационарлық вагон аударғыш, ұсақтау-фрезерлеу машинасы, жылжымалы
таспалы конвейер, бункер, бұрғылау ұнтақтау машинасы, желдету жүйесі,
қиғаш таспалы конвейер, жылжымалы қалақшалы фидер, қап сүзгісі, диірмен
желдеткіші, шлюз фидері, инерциялық көлбеу, конустық уатқыштар жатады.

Ұжымдық қорғау құралдарына КЦМП дымқыл шаң жинағыш, УГ-115 және УГ-230
типті электр сүзгілері, ротоклон скрубберлері жатады.

2022 жылы осы цехтың 1 қызметкері вагон аударғыш роторының роликті
болттарын тарту және айлау жұмыстарын жүргізу кезінде зардап шекті, оған
жұмыс өндірісін қажетті деңгейде ұйымдастырылмауы себеп болды.

Цехтың жалпы жұмысшылар саны-151 адам, 2024 жылы оның 13-і тұрмыстық
жарақат алды.
\end{multicols}

{\bfseries 1-сурет. «Ұсақтау-кептіру цехы» қызметкерлерінің жарақаттары:
зақымдануды оқшаулау және жұмыс өтілі}

Бұл 1-суретте кәсіби тәжірибе мен жас факторының жарақаттану сипатына
әсерін жан-жақты бағалауға мүмкіндік беретін зақымдануды оқшаулау
тұрғысынан жұмыс өтілі мен зардап шеккендердің жасына байланысты
өндірістік жарақаттардың таралуы көрсетілген.

Жарақат алған жұмысшылардың жасы көбінесе 25 пен 43 жас аралығында
болған, бұл еңбекке жарамды жасты көрсетеді. Зардап шеккендердің
максималды жасы 42-43 жас, жұмысшылардың бұл санатында аяқтың жарақаты,
көгеру мен қолдың сынуы көбірек кездеседі.

Еңбек өтілі бойынша жарақаттар (аяқтың сынуы және соғып алуы)
жұмысшылардың барлық санаттарында тіркеледі, бұл жұмыс өтілі неғұрлым
көп болса, жарақат түрі соғұрлым аз болады деген сенімді жоққа шығарады.

Егер қызметкерлердің кәсіптері бойынша жарақаттарды талдайтын болсақ,
осы өндірістік цехта 29-31 жас аралығындағы электр жабдықтарын жөндеу
және техникалық қызмет көрсетуші 3-6 жыл аралығында еңбек өтілі бар
электромонтерлар қолдың сынуы жарақатын алған; 1-3 жыл жұмыс өтілі бар
22-31 жас аралығындағы мөлшерлеу аппаратшысы аяғын соғып алған; 25 және
32 жастағы электрдәнекерлеушілер аяқ сынуы және соғып алу жарақаттарын
алғанын байқаймыз.

Осылайша деректерді талдау көрсеткендей, зерттеу жүргізілген цехта жұмыс
өндірісін тиісті деңгейде ұйымдастырылмауы себепті биіктен құлау кезінде
дененің жұмсақ тіндерінің көгеруі мен жаралануымен бір өндірістік және
13 тұрмыстық жарақат орын алды. Ең жиі кездесетін жарақат түрлері --
аяқтың жоғары және төменгі бөліктерінің жаралануы (соғып алу, сыну).
Зардап шеккендердің негізгі үлесі-25-43 жас аралығындағы жұмысшыларға
тиесілі.

\emph{«Агломерация»} цехында үш агломерация өндірісі желісі жұмыс
істейді. Агломерация процесі фосфорит ұнтағын AKM-7-312 агломерация
машинасында 1623 К дейінгі температурада қатты отын ретінде ұсақталған
коксты пайдаланып агломерациялауды қамтиды. Алынған агломерат
салқындатылып, қажетті мөлшерге дейін ұсақталады және сұрыпталғаннан
кейін пеш цехының дайындау бөліміне жіберіледі.

Цехта негізгі технологиялық жабдықтар, мысалы, науа таспалы конвейер,
конвейерлік таразы, жылжымалы қайтымды таспалы конвейер, силос,
электрлік дірілдейтін қоректендіруші, көлденең таспалы конвейер,
үлестіргіш, көлбеу таспалы конвейер, бункер, өлшеу құрылғысы, қайтару
үлестіргіш, барабанды араластырғыш, күйдіру машинасы, 26 вакуумдық
камерасы бар газ коллекторы, шихта түйіршіктегіші, шаттлдық заряд
таратушы, барабанды заряд беруші, диірмен желдеткіші, бір орамды
ұсатқыш, қырғыш конвейері қолданылады.

Өндіріспен байланысты жазатайым оқиғалардан зардап шеккендерді қоса
алғанда, кәсіптік тәуекелдерге ұшыраған адамдар (6 адам):

- түйіршіктегіш операторы (бейорганикалық шаң, фосфин, жұмыскердің
қозғалыс кезінде құлауы, өндірістік жабдықтың үшкір бұрыштары бар
элементтерінің әсеріне ұшырауы, жұмыскердің қозғалыс кезінде жабдықтың
бетінде жоғары және төмен температуралардың әсеріне ұшырауы, электр
тогының соғуы) (2023 жылғы жазатайым оқиға);

- қашықтан басқару операторы (бейорганикалық шаң, фосфин, шу) (2023
жылғы апат); - жөндеуші (жарықтың жеткіліксіздігі, шу, фосфин,
бейорганикалық шаң, жұмысшының қозғалыс кезінде құлауы, өткір бұрыштары
бар өндірістік жабдықтың құрылымдық элементтеріне әсер ету) (2023 жылы 2
апат);

- электр және газ дәнекерлеушісі (бейорганикалық шаң, шу, химиялық
факторлар, өткір бұрыштары бар өндірістік жабдықтың құрылымдық
элементтеріне әсер ету) (2023 жылғы апат);

- газды тазарту операторы (шу, химиялық факторлар, агломераттан шаңның
шығарылуы, заряд, жабдықтың, механизмдердің, машиналардың қозғалатын
және айналмалы бөліктеріне әсер ету) (2023 жылғы апат);

- ыстық агломератты ұсатқыш (шу, агломерат шаңының бөлінуі, шихта,
фосфин, өткір бұрыштары бар өндірістік жабдықтың құрылымдық
элементтерінің әсері, өрт немесе жарылыс қаупі (үйкеліс немесе жоғары
қысым).

Осы цехтың 6 жұмыскері 2021-2024 жылдары өндірістегі жазатайым
оқиғалардан зардап шекті, олардың себептері жұмыс өндірісін
қанағаттанарлықсыз ұйымдастыру, еңбек қауіпсіздігі және еңбекті қорғау
ережелерін бұзу, жұмыс басшысы тарапынан бақылаудың әлсіреуі болды.

2021 және 2024 жылдар аралығында осы цехтағы 6 жұмыскер жұмыстың нашар
ұйымдастырылуынан, еңбек қауіпсіздігі және еңбекті қорғау ережелерін
бұзудан және жұмыс басшысының бақылауының әлсіреуінен туындаған жұмысқа
байланысты жазатайым оқиғалардан зардап шекті.

1. ПT-616 (16) диск беру құрылғысының позициясында +19,0 метрлік белгіде
жұмыс істеу кезінде цехтың түйіршіктеу машинасының операторы науада
тұрып қалған металды алып тастауға ниеттеніп, сол қолының 3-ші және 4-ші
саусақтарының тырнақ фалангаларын жарақаттық ампутациядан алды.

Оқиға түрі: қозғалатын, ұшатын немесе айналатын заттар мен бөлшектердің
соққысы. Оқиға себептері: жұмыскердің өрескел немқұрайлылығы;

2. Цехтың қашықтан басқару пультінің операторы бетондау бөлімінің
басқару пультіне бара жатып, шалбарын гипсокартон тіреуішіне қысып алып,
құлады. Құлау салдарынан оң жақ тізе буынының жұмсақ тіндерінің жарақаты
болды. Оқиға түрі жәбірленушінің құлауы болды. Оқиғаның себебі: жұмыс
процедураларының дұрыс ұйымдастырылмауы;

3. ПT-630 позициясында +0,0 м белгісінде таспалы конвейерге арналған
резеңке тығыздағышты дайындау кезінде, 2-ші агломашинасының /күйдіру/
өткелінде, жөндеуші механик оң жақ санына пышақ жарақатын алды. Оқиғаның
түрі- өткір затпен кесілген. Оқиғаның себебі жұмыскердің өрескел
немқұрайлылығы және жұмыс басшысының бақылауының әлсіреуі болды.

4. Электр қозғалтқышы қосылған кезде SB-860/2 араластырғышындағы беріліс
қорабының білігіндегі жұмсақ муфтаны ауыстыру кезінде электр газымен
дәнекерлеуші мен жөндеуші әртүрлі жарақаттар
алды. Оқиға қозғалатын, ұшатын немесе айналатын заттар мен бөлшектердің
соққысынан болды. Оқиғаның себебі жұмыстың нашар ұйымдастырылуы болды.

Газ тазарту операторы жаяу жүргіншілер жолымен жүріп бара жатып сүрініп
құлады. Құлау салдарынан ол бел-сегізкөз омыртқасынан жарақат алды. Апат
құлау деп жіктелді. Апаттың себептері жұмысты нашар ұйымдастыру және
еңбек қауіпсіздігі мен еңбекті қорғау ережелерін бұзу болды.

Цехтағы жұмысшылардың жалпы саны 356 адамды құрайды, оның ішінде 16
жұмысшы тұрмыстық жарақат алған.

{\bfseries 2-сурет. «Сары фосфор алу үшін фосфат агломератын өндіру»
цехындағы жұмыскерлер арасындағы жарақаттар: жарақаттардың орналасқан
жері және жұмыс өтілі}

Осы диаграмманы талдай отырып, цехта жұмыскерлердің көпшілігі аяқ пен
қол жарақаттарын (сынуы, жарақат алу/ушиб/) алғанын атап өтуге болады.
Мамандығы бойынша бөліп қарастырсақ, 4 жылдық тәжірибесі бар 31 жастағы
салқындату операторының аяғы сынған; 8 жылдық тәжірибесі бар 29 жастағы
электр газ дәнекерлеушісінің жағы сынған; 9 жылдық тәжірибесі бар 29
жастағы қашықтан басқару операторының қолы сынған; 4 жылдық тәжірибесі
бар 29 жастағы конвейер операторының басы контузия алған; ал 8 және 12
жылдық тәжірибесі бар 37 жастағы газ тазарту операторларының жамбасы
сынған және арқасынан жарақат алған.

Осылайша, деректерді талдау зерттеліп отырған цехта 5 жазатайым оқиға
және 6 жарақат алған жұмысшы, сондай-ақ тұрмыстық жарақаттың 16 жағдайы
анықталғанын көрсетті. Ең көп таралған жарақаттар жоғарғы және төменгі
аяқ-қолдарда болды (көгерулер, сынықтар). Жарақат алған жұмысшылардың
көпшілігі 29-37 жас аралығында болды.

Осылайша, деректерді талдау зерттелетін цехта 5 жазатайым оқиға мен 6
зардап шеккен жұмыскердің, сондай-ақ 16 тұрмыстық жарақаттың
анықталғанын көрсетті. Ең жиі кездесетіні-жоғарғы және төменгі
аяқтардағы жарақаттар (жарақат алуы, сынықтар). Зардап шеккендердің
негізгі үлесі 29-37 жас аралығындағы жұмыскерлерге тиесілі.

\emph{«Сары фосфор өндіретін»} цех сегіз шихталау станциясынан,
сондай-ақ фосфор өндіретін төрт технологиялық желіден тұрады. Фосфор
өндірісінің әрбір технологиялық желісіне фосфорлы электр пеші және пеш
газын шаңнан және фосфор конденсациясынан тазартатын екі параллель
жүйесі кіреді. Фосфор алудың барлық желілері РКЗ - 80Ф-И1 фосфор
пештерімен және ЭВТ-2-5,5-20Ф-01 электр сүзгілерімен жабдықталған.
Өндіріс әдісі-электротермиялық, руднотермиялық пештерде кремний диоксиді
болған кезде фосфаттарды көміртегімен тотықсыздандыру.

Цехта көлбеу науалық таспалы конвейер, көлденең науалық кері айналмалы
таспалы конвейер, кран-арқалық, пеш бункері, фосфор қабылдағышы, фосфор
конденсаторы, ортадан тепкіш сорғы, секторлық ысырма, фосфор пеші, тік
фосфорлы электр сүзгісі, суасты сорғысы, секторлық ысырма, сақтандырғыш
гидроысырма сияқты негізгі технологиялық жабдықтар қолданылады.

Ұжымдық қорғаныс құралдарына ЦН-15 циклоны, ФРКДИ-1100 қап сүзгісі,
СКЦН-34 циклоны және ФРКИ-360 қап сүзгілері кіреді. Барлығы 454 жұмысшы
жұмылдырылды, оның 22-сі тұрмыстық жарақат алды.

{\bfseries 3-сурет. Сары фосфор өндірісі цехындағы жұмысшылар арасындағы
жарақаттар: жарақаттың орны және еңбек өтілі}

Ұсынылған диаграмма зардап шеккендердің жасы мен еңбек өтілін ескере
отырып, жарақат түрлері бойынша цех жұмысшылары арасындағы өндірістік
жарақаттың құрылымын көрсетеді, бұл жарақаттардың сипатына кәсіби және
демографиялық факторлардың әсерін бағалауға мүмкіндік береді.

Талдау көрсеткендей, сары фосфор цехындағы жарақаттар негізінен 26-49
жас аралығындағы жұмысшыларда тіркеледі, ал ең үлкен жас мәндері зардап
шеккендердің жасы 45-49 жасқа жететін төменгі аяқтардың көгеруі мен
дислокациясы сияқты зақымдануларға тән. Бұл өндірістік операцияларға
белсенді қатысатын орта және жоғары еңбекке қабілетті жастағы жұмысшылар
арасында жоғары тәуекел деңгейін көрсетеді.

Еңбек өтілі бойынша жарақаттар ең аз жұмыс тәжірибесі бар жұмысшыларда
(2-4 жыл) және 20 жылға дейін немесе одан да көп жұмыс өтілі бар
жұмысшыларда тіркеледі. Еңбек өтілі бойынша ең маңызды көрсеткіштер
аяқтың шығуы мен сынуы, сондай-ақ қолдың кесілуі мен жарақаттануы
кезінде тіркелді, бұл қауіпті өндірістік факторлардың ұзақ мерзімді
әсерінің, сондай-ақ кәсіптік тәртіпке байланысты назар мен сақтық
деңгейінің төмендеуін көрсетуі мүмкін.

Әсіресе, тірек-қимыл аппаратының зақымдануының - аяқтың, қабырғаның
сынуы, сондай-ақ аяқтың шығуы мен көгеруі көп таралғанын атап өтуге
болады. Бұл жұмыс орындарын ұйымдастыруға, өндірістік жабдықтың
жай-күйіне, тауарлардың қозғалысына және цех ішіндегі қозғалыс
қауіпсіздігіне байланысты ықтимал проблемаларды көрсетеді.

Тәжірибесі аз (2-4 жыл) жұмысшыларда қол жарақаттары мен аяқтың сынуы
сияқты жарақаттардың салыстырмалы түрде аз үлесі басшылықтың қатаң
бақылауымен және бейімделу кезеңінде еңбекті қорғау талаптарын сақтауға
көбірек көңіл бөлумен байланысты, сонымен қатар тәжірибесі жеткіліксіз
жұмысшылардың белгілі бір сақтығымен болуы мүмкін.

Жалпы, талдау нәтижелері цехтағы өндірістік жарақаттанудың жұмысшылардың
жасына да, еңбек өтіліне де тәуелділігі бар деген қорытынды жасауға
мүмкіндік береді. Бұл мақсатты алдын алу шараларын енгізу, соның ішінде
тәжірибелі жұмысшыларға қауіпсіз жұмыс әдістерін оқытуды күшейту, кәсіби
тәуекелдерді үнемі бағалау, жұмыс орындарының эргономикасын жақсарту
және еңбек қауіпсіздігі мәдениетін арттырудың қажеттілігін көрсетеді.

Осы цехтағы жұмысшылардың жарақаттарын кәсібі бойынша талдай отырып,
28-49 жас аралығындағы, 2-24 жыл жұмыс өтілі бар сары фосфор өндірісінің
операторлары жарақаттанғанын, соның ішінде орнынан шығу, контузия,
аяқтың сынуы және қабырғасының сынғанын атап өтуге болады.

Кәсіптік қауіп-қатерге ұшырағандар, соның ішінде өндірістегі жазатайым
оқиғалардан зардап шеккендер (8 адам):

- слесарь-сантехник (бейорганикалық кварцит, кокс, фосфорит шаңы, шу,
биіктіктен құлау, құлау, опырылу, жұмысшыға заттардың құлауы, өткір
бұрыштары бар өндірістік жабдықтардың құрылымдық элементтерінің әсері,
электр тогының соғуы) (2022 жылғы жазатайым оқиға);

- сары фосфор (түйіршіктеу) өндіру аппаратшысы (шу, фосфор оксиді,
фосфин, ақ фосфор, биіктіктен құлау, қозғалу кезінде құлау, опырылу,
жұмыскерге заттардың құлауы, ғимараттардың, құрылыстардың және олардың
элементтерінің құлауы, қирауы, жабдықтың қозғалатын және айналатын
бөліктерінің әсері, электр тогының соғуы) (жазатайым оқиға 2024 жылы);

- сары фосфор өндірісінің аппаратшысы (шихталаушы) (жарылыс, пеш газы,
изотоптық датчиктер, жоғары температура T 1400 \tsp{o}C)
(2024 жылғы жазатайым оқиға);

- сары фосфор аппаратшысы (аға) (химиялық факторлар, өткір жиектер,
қабыршақтар, құрал-сайманның кедір-бұдырлы беті, қозғалатын шикізат,
материалдар, жоғары шаңдану және газдану, электр тогынан зақымдану)
(2023 жылы және 2024 жылы жазатайым оқиға);

- электр жабдықтарына қызмет көрсету жөніндегі электр монтері
(органикалық емес шаң, электр тогымен зақымдану, өрт немесе жарылыс
қаупі, биіктіктен құлау, құлау, ғимараттар мен олардың элементтерінің
қирауы, жабдықтардың, механизмдердің, машиналардың қозғалатын және
айналатын бөліктерінің әсері) (2024 жылғы жазатайым оқиға);

- сары фосфор өндірісінің аппаратшысы (ферроқорытпаларды ағызу бойынша)
(шашыраулар, жоғары температура T 1400 \tsp{o}C, электр
тогымен зақымдану, көміртек тотығымен улану, өрт немесе жарылыс қаупі);

- сары фосфор (конденсация) өндірісінің аппаратшысы (шу, фосфор оксиді,
электр тогымен зақымдану, өрт немесе жарылыс қаупі, жабдықтың,
механизмдердің, машиналардың қозғалатын және айналатын бөліктерінің
әсері, үшкір бұрыштары, жиектері, қабыршақтары және тегіс емес беттері
бар өндірістік жабдық конструкциясы элементтерінің әсері, сондай-ақ
жұмыскердің қозғалысы кезінде жабдық бетінің жоғары және төмен
температурасының әсері) (жазатайым оқиға 2023 жылы);

- сары фосфор өндіру аппаратшысы (шаң-газ ұстайтын қондырғыларға қызмет
көрсету бойынша) (шу, фосфин, ақ фосфор, үшкір бұрыштары, жиектері,
қабыршақтары және тегіс емес беттері бар өндірістік жабдық конструкциясы
элементтерінің әсері, сондай-ақ жұмыскердің қозғалысы кезінде жабдық
бетінің жоғары және төмен температурасының әсері, жабдықтың,
механизмдердің, машиналардың қозғалатын және айналатын бөліктерінің
әсері, өрт немесе жарылыс қаупі).

2021-2024 жылдары өндірістегі жазатайым оқиғалардан осы цехтың 9
қызметкері зардап шекті, себептердің бірі - жұмыс өндірісінің
қанағаттанарлықсыз ұйымдастырылуы болып табылады.

1. Осы цехтың сары фосфор өндіру аппаратшысы Батыс Қазақстан облысы
Ақсай қаласының «Қазақстан» станциясында іссапарда болған кезде теміржол
құрамының астына түсіп қалған. Диагноз: оң аяқтың тізе буыны деңгейінде
травматикалық кесілуі. Оқиға түрі - темір жол көлігі оқиғасы. Оқиға
себептері - қызметкердің өрескел абайсыздығы және жұмыс өндірісінің
қанағаттанарлықсыз ұйымдастырылуы;

2. Слесарь-сантехник цехтың әкімшілік-тұрмыстық корпусының
жертөлесіндегі кәріз құбырындағы ағып кетуді жою жұмыстарын жүргізген
кезде қоқыстан құлады. Құлау салдарынан жабық бас сүйек жарақаты алған.
Мидың соғылуы. Жедел субдуральды гематома, оң жақта. Оқиға түрі - зардап
шегушінің құлауы. Оқиға себептері - қызметкердің өрескел абайсыздығы
және жұмыс өндірісінің қанағаттанарлықсыз ұйымдастырылуы;

3. Сары фосфор өндірісінің (кәдеге жарату) аппаратшысы «Г» ауысымында
5-61-6 позицияда сорғы сальнигін толтыру кезінде конденсация
бөлімшесінде жұмыс істегенде 0,8 м белгі көз қабығының жеңіл дәрежелі
шамалы күйігін алды. Оқиға түрі - зиянды және қауіпті өндірістік
факторлар мен заттардың әсері. Оқиға себептері - қызметкердің өрескел
абайсыздығы және жұмыс өндірісінің қанағаттанарлықсыз ұйымдастырылуы;

4. Цехтың сары фосфор өндірісінің аға аппаратшысы «Г» ауысымында жұмыс
істеп, № 6 пештің секторлық қақпағының 9 тиеу ағынының жанынан + 22,4 м
белгісінде ессіз күйде табылды. Алғашқы медициналық көмек көрсетілгеннен
кейін қалалық ауруханаға жеткізілді. Оқиға түрі-зиянды және қауіпті
өндірістік факторлар мен заттардың әсері. Оқиға себептері-жұмыс
өндірісінің қанағаттанарлықсыз ұйымдастырылуы;

5. Осы жұмыстарды жүргізу кезінде сары фосфор өндірісінің аппаратшысы
НУВ-12, қызмет көрсету алаңының баспалдақының шетіне тайып түсіп, 16,4
метр белгі алаңына арқасымен құлады. Құлау нәтижесінде сары фосфор
өндірісінің аппаратшысы оң жақтағы VI-VII қабырғаларының жабық сынығын
алды. Оқиға түрі-зардап шегушінің құлауы. Оқиғаның себептері-зардап
шегушінің өрескел абайсыздығы;

6. Цехтың пеш бөлімшесінде технологиялық жабдықты тексеру (аралау)
барысында +43,0 метрлік белгісі бар баспалдақпен +40,0 метрлік белгімен
түсіп, цехтың сары фосфор (шихталаушы) өндірісінің аппаратшысы құлады.
Диагноз: сол жақ сан сүйегінің үлкен ұршығының жабық үзік сынығы. Оқиға
түрі - зардап шегушінің құлауы. Оқиғаның себептері - зардап шегушінің
өрескел абайсыздығы;

7. Цехтың сары фосфор өндіру аппаратшысы асханаға бара жатып,
2-түйіршіктеу бөлімінде оң аяқты +0,00 м белгісінде қайырып алады.
Диагноз: Оң жіліншік буынының байламдарының созылуы. Оқиға түрі-зардап
шегушінің құлауы. Оқиғаның себептері- зардап шегушінің өрескел
абайсыздығы;

8. Цехтың пеш бөлімшесінде 7,8 метр белгідегі №6 пештің ферроледкасын
жабу бойынша жұмыстарды жүргізу кезінде ферраның шашырауы сары фосфор
(аға) өндіру аппаратшысының оң көзіне түсті. Диагноз: оң көздің
конъюнктивасы мен қабығының термохимиялық күйігі. Оқиға түрі-зиянды және
қауіпті өндірістік факторлар мен заттардың әсері. Оқиғаның
себептері-жәбірленушінің өрескел абайсыздығы;

9. Электродты масса учаскесінің аға шебері. Диагноз: оң жақтағы өкше
сүйегінің жабық сынуы. Оқиға түрі-әбірленушінің биіктіктен құлауы.
Оқиғаның себептері-жәбірленушінің өрескел абайсыздығы;

10. Цехтың пеш бөлімшесінде +7,8 метр белгідегі №7 пештің пеш
трансформаторларының камерасында трансформаторды тексеру жұмыстарын
жүргізген кезде электр жабдықтарына қызмет көрсету жөніндегі
электромонтер электр тогымен зақымданды. Жағдайлар, себептер және
ауырлық дәрежесі анықталады. Алдын ала диагноз: дене бетінің 80\% күйік.
Оқиғаның түрі-электр тогының соғуы. Оқиғаның себептері-жәбірленушінің
өрескел абайсыздығы;

Деректерді талдауды қорытындылай келе, зерттелетін цехта жұмыс
өндірісінің қанағаттанарлықсыз ұйымдастырылуына және жәбірленушінің
өрескел абайсыздығына, сондай-ақ тұрмыстық жарақаттардың 22 жағдайына
байланысты өндірістік жарақаттанудың 9 жағдайы анықталғанын көрсетті. Ең
жиі кездесетіні-жоғарғы және төменгі аяқтардағы жарақаттар (көгеру,
сынықтар). Зардап шеккендердің негізгі үлесі 26-45 жас аралығындағы
жұмысшыларға тиесілі.

\emph{«Сары фосфорды өңдеу»} цехында тауарлық фосфорды алу
тұндырғыштарда шикі фосфорды тұндыру, теміржол цистерналарында,
контейнерлерде және бөшкелерде фосфорды жіберу, теміржол цистерналарын
жуу әдісімен жүзеге асырылады. Сондай-ақ, техникалық фосфорды
органикалық қоспалардан және мышьяк қоспаларынан тазарту жүргізіледі.
Техникалық фосфорды тазарту дәрежесі тұтынушылардың талабы бойынша
жасалады.

Бұл цехта тұндырғыш қоюландырғыш, буферлік сыйымдылық,
бейтараптандырғыш, сары фосфор тұндырғышы, батырмалы сорғы, ауа үрлегіш,
вакуумдық сорғы, инверторлы дәнекерлеу аппараты, ФТН 1253К центрифуга
сияқты негізгі технологиялық жабдықтар қолданылады.

Ұжымдық қорғау құралдарының ішінде: циклондық қондырғы, ГПФ7-9 электр
сүзгілері, циклофильтр.

Өндіріске байланысты жазатайым оқиғалардан зардап шеккендерді қоса
алғанда, кәсіби тәуекелдерге ұшырағандар (3 адам):

- сары фосфор өндірісінің аппаратшысы (айналмалы және қозғалмалы
механизмдердің, қысыммен жұмыс істейтін аппаратуралардың, фосфордың,
фосшламның, сутегі фторидінің, фосфиннің уытты әсерінің болуы) (2021,
2022 және 2023 жылдардағы 3 жазатайым оқиға);

- электр жабдықтарын жөндеу және техникалық қызмет көрсету жөніндегі
электромонтер (айналмалы және қозғалмалы механизмдердің болуы,
биіктіктен құлау, фосфиннің, ақ фосфордың уытты әсері, электр тогының
соғуы, өрт немесе жарылыс қаупі) (2023 жылғы жазатайым оқиға);

- ілмектеуші (биіктіктен құлау, кальцийленген соданың бейорганикалық
тозаңы, электр тогының зақымдануы, өрт немесе жарылыс қаупі,
жабдықтардың, механизмдердің, машиналардың қозғалатын және айналатын
бөліктерінің әсері) (2023 жылғы жазатайым оқиға).

2021-2024 жылдары аталған цехтың 5 қызметкері өндірістегі жазатайым
оқиғалардан зардап шекті, олардың себептері жұмыс өндірісінің
қанағаттанарлықсыз ұйымдастырылуы, жеке қорғану құралдарын қолданбаудан
болды.

1. Аталған цехтың сары фосфорды өндіру аппаратшысын Е-1670/2
позициясының сыйымдылығынан Ду = 25 мм резеңке жеңін шығару кезінде
құрамында фосфоры бар ағындар шашыраған. Диагноз: Екі аяқ-қолдың әртүрлі
дәрежедегі фосфор шламымен термохимиялық күйік. Оқиға түрі - зиянды
өндірістік факторлардың әсері. Оқиға себептерi - жұмыс өндiрiсiн
қанағаттанарлықсыз ұйымдастыру және жеке қорғану құралдарын қолданбау.

2. П-6 тазартылған фосфор позасының қабылдағышында өлшеуішпен фосфорды
өлшеу кезінде цехтың сары фосфор өндірісінің аппаратшысы фосфор II-ІІІА
буымен бет, мойын және екі білек, денесінде S\textasciitilde{} 4\%
термохимиялық күйік алды. Оқиғаның түрі-зиянды өндірістік факторлардың
әсері. Оқиғаның себептері - жұмыс өндірісін қанағаттанарлықсыз
ұйымдастыру және жеке қорғаныс құралдарын қолданбау;

3. Цехтың электр жабдықтарын жөндеу және қызмет көрсету жөніндегі электр
монтері 1309/4А позициясының сорғысын жинау кезінде сол қолының
саусақтарынан күйік алған. Оқиға түрі - зиянды өндірістік факторлардың
әсері. Оқиға себептері - жұмыс өндірісін қанағаттанарлықсыз ұйымдастыру
және жеке қорғану құралдарын қолданбау;

4. Е 1331/2 позициясының қоймасынан Е 1/8 позициясының монжусына
техникалық сары фосфорды қабылдау жөніндегі жұмыстарды жүргізу кезінде
сары фосфор өндірісінің аппаратшысы фосфорды сумен лақтырды. Лақтыру
кезінде сары фосфор өндірісінің аппаратшысы қызмет көрсету алаңынан 0,0
метр шамасында 5,8 метр шамасында секірді. Биіктіктен құлау нәтижесінде
зардап шегуші сол жақ жамбастың жабық сынықтарын алды. Екі санның, кеуде
қуысының, сол қолдың термиялық, химиялық күйігі 10 \% /1-2-3 кезең.
Оқиға түрі - зиянды өндірістік факторлардың әсері. Оқиға себептері -
жұмыс өндірісін қанағаттанарлықсыз ұйымдастыру және жеке қорғану
құралдарын қолданбау;

5. Сары фосфор қоймасының аумағында осы цехтың ілмектеуші дененің
әртүрлі бөліктерінде химиялық күйік алды. Оқиғаның түрі -- химиялық
күйік. Оқиғаның себептері -- жұмыс өндірісін қанағаттанарлықсыз
ұйымдастыру.

Цехтағы жұмысшылардың жалпы саны 145 адамды құрайды, оның ішінде 7
жұмысшы тұрмыстық жарақат алған.

{\bfseries 2-сурет. «Фосфорды өңдеу, бейтараптандыру және жағу» цехындағы
жұмыскерлер арасындағы жарақаттар: зақымдануды оқшаулау және жұмыс
өтілі}

Осы диаграммаларды талдай отырып, цехтағы жұмысшылардың көбісінде
қолдары мен аяқтарына жарақат алған (сыну, ұрып алу) байқауға болады.
Егер 41 мен 27 жастағы слесарь-жөндеуші мамандығы бойынша егжей-тегжейлі
айтсақ, еңбек өтілі 18 жыл және 3 жыл аяқ-қолдары жарақат алған. Ал 21
жастағы электр газбен дәнекерлеуші, еңбек өтілі 1 жыл ғана аяқтары
сынған, электр жабдықтарын жөндеу және қызмет көрсету жөніндегі
электромонтер - еңбек қызметі нәтижесінде 2 жыл өтілі бар 22 жастағы
жұмыскер аяғынан жарақат алған болса, 31 жастағы еңбек өтілі 2 жыл
ілмектеуші сол қолының химиялық күйігін алды.

Деректерді талдау негізінде зерттелетін цехта өндірістік жарақаттанудың
5 жазатайым оқиғасы және тұрмыстық жарақаттардың 7 жағдайы анықталғанын
көрсетті. Ең жиі кездесетіні жарақаттардың түрі-аяқ-қолдардың
жарақаттары (ұрып алу, соғып алу, сынуы, орынынан шығып кетуі). Зардап
шеккендердің негізгі үлесі 21-27 жас аралығындағы жұмысшыларға тиесілі.

\emph{«Азот-оттегі өндірісі»} цехында сары фосфор өндірісінің
технологиялық сызбасының құрамына кіреді және азот газын, сығылған
құрғатылған және кептірілмеген ауаны, техникалық оттегі газын өндіруге
және зауытқа беруге қызмет етеді. Сұйық азот пен сұйық техникалық
оттегін жағына жібереді. «Азот-оттегі өндірісі» ауаны төмен
температурада ректификациялау әдісімен жүзеге асырылады.

Технологиялық сызба төмен қысымды цикл бойынша ауаны кешенді тазарту
(АКТ) блогын, турбодетандерлік агрегатты (ДТГ411) қолдана отырып, ағызу
ағынында газ үрлегішпен салынған. Негізгі бөлу аппараты бір реттік
ректификация сызбасы бойынша салынған.

Бұл цехта ауаны бөлу блогы, АКТ-30 сұйытылған газ сорғысы, «К-500»
орталық тепкіш компрессор, «ДТГ-14/0,3» турбодетандер, ауаны кептіру
блогы, ауа компрессоры сияқты негізгі технологиялық жабдықтар
қолданылады.

{\bfseries 5-сурет. «Азот-оттегі өндіріс» цехының слесарь-жөндеушісінің
қолын соғып алу: зақымдануды оқшаулау және жұмыс өтілі}

\begin{multicols}{2}
Өндіріске байланысты кәсіптік тәуекелдерге ұшырағандар (3 адам):

- ауа бөлу аппаратшысы (шу, биіктіктен құлау, қозғалу кезінде құлау,
заттардың құлауы, құлау, жабдықтың, механизмдердің, мәшиналардың
қозғалатын және айналатын бөліктерінің әсері, электр тоғының соғуы, өрт
немесе жарылыс қаупі);

- газды кептіру аппаратшысы (шу, биіктіктен құлау, қозғалу кезінде
құлау, заттардың құлауы, құлау, жабдықтың, механизмдердің,
мәшиналардың қозғалатын және айналатын бөліктерінің әсері, электр
тоғының соғуы, өрт немесе жарылыс қаупі);

- баллондарды толтырғыш (шу, химиялық факторлар, қысымдағы ыдыстар,
биіктіктен құлау, қозғалу кезінде құлау, заттардың құлауы, құлау,
жабдықтың, механизмдердің, мәшиналардың қозғалатын және айналатын
бөліктерінің әсері, электр тоғының соғуы, өрт немесе жарылыс қаупі).

2021-2024 жылдары бұл цехта жазатайым оқиғалар тіркелген жоқ. Жалпы саны
56 жұмыскер, оның ішінде 1 жұмыскер тұрмыстық жарақат алды.

Слесарь-жөндеуші, 26 жас, 5 жыл еңбек өтілі, қолын соғып алу.

Осы цех бойынша соңғы 2020-2025 жылдарда өндірістік жарақаттану
анықталған жоқ, бірақ тұрмыстық жарақаттың бір жағдайы орын алды. Осы
кәсіпорында 5 жылдық жұмыс өтілі бар 26 жастағы жұмыскер
(слесарь-жөндеуші).

{\bfseries Қорытынды.} 2024 жылғы жалпыланған деректер фосфор өнеркәсібі
кәсіпорындарында жарақаттанудың 65 жағдайы тіркелгенін көрсетеді, оның
6-ы өндірістік, 59-ы тұрмыстық. Өндірістік жарақаттану көрсеткіштері
жазатайым оқиғаларды есепке алу мен тіркеудің толық болмауына байланысты
саладағы қауіпсіздік жағдайы мен тәуекелдерді басқару деңгейінің дұрыс
бағалануын қамтамасыз етпейді. Тіркелген өндірістік жарақаттардың төмен
деңгейі кәсіптік тәуекелдерді объективті бағалау мүмкіндіктерін шектейді
және еңбекті қорғауды басқарудың тиімділігін төмендетеді. Алынған
деректер еңбек қауіпсіздігін нормативтік-құқықтық және ұйымдастырушылық
қамтамасыз етуді жетілдіру және кәсіптік тәуекелдерді басқаруға
тәуекел-тәсілді енгізу қажеттілігін негіздейді.
\end{multicols}

\begin{center}
{\bfseries Әдебиеттер}
\end{center}

\begin{refs}
1. Қазақстан Республикасының Еңбек Кодексі.2015 жылғы 23 қарашадағы
№414-V ҚРЗ.1 қаулы, 61 тармақша. Электрондық ресурс.
-URL:\href{https://adilet.zan.kz/rus/docs/K1500000414}{adilet.zan.kz/rus/docs/K1500000414}.
\url{https://adilet.zan.kz/kaz/docs/Z1400000243.-Қаралған} күні:
25.12.2025.

2. Шаяхметова, Л., Курманов, А., \& Айтимова, Ш. (2024).
Совершенствование охраны труда в Республике Казахстан: оценка
профессиональных рисков и стратегии страхования от несчастных случаев на
производстве// Bulletin the National academy of sciences of the Republic
of Kazakhstan.-
Vol.5(411).-P.401-418.\href{https://doi.org/10.32014/2024.2518-1467.849}{DOI
10.32014/2024.2518-1467.849}.

3. Afidah A., Khamisah Awang L., Helmy S., Syed Sharizman Syed A.,
Fredie R., Mohd Rohaizat H. Prevalence of occupational injury and
determination of safety climate in small scale manufacturing industry: A
cross-sectional study //Annals of Medicine \&
Surgery\href{file:///C:/Users/admin/Downloads/69.\%20-\%20September\%202021}{.
- 2021.}-Vol.69: 102699.
\href{https://doi.org/10.1016/j.amsu.2021.102699}{DOI
10.1016/j.amsu.2021.102699}.

4. Ердесов Н.Ж., Сраубаев Е.Н., Cерик Б. Производственный травматизм и
профессиональная заболеваемость в Республике Казахстан //Экология и
медицина.2020. - № 4. - С.38- 45.

5. Ghahramani A., Amirbahmani A. A Study of the Causes of Occupational
Accidents in Manufacturing
Companies//\href{https://www.researchgate.net/journal/Archives-of-Trauma-Research-2251-9599?_tp=eyJjb250ZXh0Ijp7ImZpcnN0UGFnZSI6InB1YmxpY2F0aW9uIiwicGFnZSI6InB1YmxpY2F0aW9uIn19}{Archives
of Trauma
Research}.-2021.-Vol.10(2):64.\href{https://doi.org/10.4103/atr.atr_56_20}{DOI
10.4103/atr.atr\_56\_20}.

6. ВОЗ оценивает стоимость выполнения глобальных задач, связанных со
здоровьем, к 2030 г. URL:
https://www.who.int/ru/news/item/17-07-2017-who-estimates-cost-of-reaching-global-health-targets-by-2030.-
\href{https://adilet.zan.kz/kaz/docs/Z1400000243.-Қаралған}{Қаралған}
күні: 25.12.2025.

7. Еңбек қауіпсіздігі мен гигиенасына жәрдемдесетін негіздер туралы
конвенцияны (187-конвенция) ратификациялау туралы Қазақстан
Республикасының Заңы 2014 жылғы 20 қазандағы № 243-V ҚРЗ. -URL:
https://adilet.zan.kz/kaz/docs/Z1400000243.-
\href{https://adilet.zan.kz/kaz/docs/Z1400000243.-Қаралған}{Қаралған}
күні: күні: 25.12.2025.

8. Қазақстан Республикасының 2024-2030 жылдарға арналған қауіпсіз еңбек
тұжырымдамасын бекіту туралы Қазақстан Республикасы Үкіметінің 2023
жылғы 26 желтоқсандағы № 1182 қаулысы.
-URL:\href{https://adilet.zan.kz/rus/docs/P2200000245}{adilet.zan.kz/rus/docs/P2200000245}.
\href{https://adilet.zan.kz/kaz/docs/Z1400000243.-Қаралған}{Қаралған}
күні: күні: 25.12.2025.

9. Rosa Maria P., Angel Tejada P., Maria Pilar S.//How to prevent 3
million deaths worldwide: a systematic review of occupational accident
research factors and cost-based approach. //EUPHA.
-2025.-Vol.35(1).-P.91-100.
\href{https://doi.org/10.1093/eurpub/ckae197}{DOI
10.1093/eurpub/ckae197}.

10. Bravo G., Viviani C., Lavallière M., Arezes P., et.all. Do older
workers suffer more workplace injuries? A systematic review.//
International Journal of Occupational Safety and Ergonomics.-
2022. -Vol.28(1)\href{https://doi.org/10.1080/10803548.2020.1763609}{. -
P.398- 427.DOI 10.1080/10803548.2020.1763609}.

11. Бакирова А.Б., Карамова Л.М., Каримова Л.К., Власова Н.В., Шаповал
И.В., Башарова Р.Г. Современные проблемы производственного травматизма
со смертельным исходом. // Гигиены труда. - 2024. -Т.1. -С.25-48. DOI
10.24412/2411-3794-2024-10102.

12. Жеглова А.В., Яцына И.В., Гаврильченко Д.С. Корпоративные программы
сохранения здоровья-основной элемент системы здоровьесбережения
работающего населения // Здравоохранение Российской Федерации. -2022.
-Т.66(5). -С.385-389.
\href{https://doi.org/10.47470/0044-197X-2022-66-5-385-389}{DOI
10.47470/0044-197X-2022-66-5-385-389}.

13. 2024 жылғы Қазақстан Республикасындағы еңбек қызметіне байланысты
жарақаттану және кәсіптік аурулар туралы//Ұлттық статистика бюросының
анықтамалығы. -2025.
-URL:\href{https://stat.gov.kz/ru/news/o-travmatizme-svyazannom-s-trudovoy-deyatelnostyu/?sphrase_id=1588311}{stat.gov.kz/ru/news/o-travmatizme-svyazannom-s-trudovoy-deyatelnostyu/?sphrase\_id=1588311}.
- \href{https://adilet.zan.kz/kaz/docs/Z1400000243.-Қаралған}{Қаралған}
күні: күні: 28.12.2025.
\end{refs}

\begin{center}
{\bfseries References}
\end{center}

\begin{refs}
1. Қazaқstan Respublikasynyң Eңbek Kodeksі.2015 zhylғy 23 қarashadaғy
№414-V ҚRZ.1 қauly, 61 tarmaқsha. Jelektrondyқ resurs.
-URL:adilet.zan.kz/rus/docs/K1500000414.
https://adilet.zan.kz/kaz/docs/Z1400000243.-Қaralғan kүnі: 25.12.2025.
{[}in Kazakh{]}

2. Shajahmetova, L., Kurmanov, A., \& Ajtimova, Sh. (2024).
Sovershenstvovanie ohrany truda v Respublike Kazahstan: ocenka
professional' nyh riskov i strategii strahovanija ot
neschastnyh sluchaev na proizvodstve// Bulletin the National academy of
sciences of the Republic of Kazakhstan. - Vol.5(411). -P.401-418.DOI
10.32014/2024.2518-1467.849. {[}in Russian{]}

3. Afidah A., Khamisah Awang L., Helmy S., Syed Sharizman Syed A.,
Fredie R., Mohd Rohaizat H. Prevalence of occupational injury and
determination of safety climate in small scale manufacturing industry: A
cross-sectional study //Annals of Medicine \&
Surgery\href{file:///C:/Users/admin/Downloads/69.\%20-\%20September\%202021}{.
- 2021.}-Vol.69: 102699.
\href{https://doi.org/10.1016/j.amsu.2021.102699}{DOI
10.1016/j.amsu.2021.102699}.

4. Erdesov N.Zh., Sraubaev E.N., Cerik B. Proizvodstvennyj travmatizm i
professional' naja zabolevaemost'{} v
Respublike Kazahstan //Jekologija i medicina.2020. - № 4. - S.38- 45.
{[}in Russian{]}

5. Ghahramani A.,Amirbahmani A. А Study of the Causes of Occupational
Accidents in Manufacturing
Companies//\href{https://www.researchgate.net/journal/Archives-of-Trauma-Research-2251-9599?_tp=eyJjb250ZXh0Ijp7ImZpcnN0UGFnZSI6InB1YmxpY2F0aW9uIiwicGFnZSI6InB1YmxpY2F0aW9uIn19}{Archives
of Trauma
Research}.-2021.-Vol.10(2):64.\href{https://doi.org/10.4103/atr.atr_56_20}{DOI
10.4103/atr.atr\_56\_20}.

6. VOZ ocenivaet stoimost'{} vypolnenija
global' nyh zadach, svjazannyh so
zdorov' em, k 2030 g. URL:
https://www.who.int/ru/news/item/17-07-2017-who-estimates-cost-of-reaching-global-health-targets-by-2030.-
Қaralғan kүnі: 25.12.2025. {[}in Russian{]}

7. Eңbek қauіpsіzdіgі men gigienasyna zhәrdemdesetіn negіzder turaly
konvencijany (187-konvencija) ratifikacijalau turaly Қazaқstan
Respublikasynyң Zaңy 2014 zhylғy 20 қazandaғy № 243-V ҚRZ. -URL:
https://adilet.zan.kz/kaz/docs/Z1400000243.- Қaralғan kүnі: kүnі:
25.12.2025. {[}in Kazakh{]}

8. Қazaқstan Respublikasynyң 2024-2030 zhyldarғa arnalғan қauіpsіz eңbek
tұzhyrymdamasyn bekіtu turaly Қazaқstan Respublikasy Үkіmetіnің 2023
zhylғy 26 zheltoқsandaғy № 1182 қaulysy.
-URL:adilet.zan.kz/rus/docs/P2200000245. Қaralғan kүnі: kүnі:
25.12.2025. {[}in Kazakh{]}

9. Rosa Maria P., Angel Tejada P., Maria Pilar S.//How to prevent 3
million deaths worldwide: a systematic review of occupational accident
research factors and cost-based approach. //EUPHA.
-2025.-Vol.35(1).-P.91-100.
\href{https://doi.org/10.1093/eurpub/ckae197}{DOI
10.1093/eurpub/ckae197}.

10. Bravo G., Viviani C., Lavallière M., Arezes P., et.all. Do older
workers suffer more workplace injuries? A systematic review.//
International Journal of Occupational Safety and Ergonomics.-
2022. -Vol.28(1)\href{https://doi.org/10.1080/10803548.2020.1763609}{. -
P.398- 427.DOI 10.1080/10803548.2020.1763609}.

11. Bakirova A.B., Karamova L.M., Karimova L.K., Vlasova N.V., Shapoval
I.V., Basharova R.G. Sovremennye problemy proizvodstvennogo travmatizma
so smertel' nym ishodom. // Gigieny truda. - 2024. -T.1.
-S.25-48. DOI 10.24412/2411-3794-2024-10102. {[}in Russian{]}

12. Zheglova A.V., Jacyna I.V., Gavril' chenko D.S.
Korporativnye programmy sohranenija zdorov' ja-osnovnoj
jelement sistemy zdorov' esberezhenija rabotajushhego
naselenija // Zdravoohranenie Rossijskoj Federacii. -2022. -T.66(5).
-S.385-389. DOI 10.47470/0044-197X-2022-66-5-385-389. {[}in Russian{]}

13. 2024 zhylғy Қazaқstan Respublikasyndaғy eңbek қyzmetіne bajlanysty
zharaқattanu zhәne kәsіptіk aurular turaly//Ұlttyқ statistika bjurosynyң
anyқtamalyғy. -2025.
-URL:stat.gov.kz/ru/news/o-travmatizme-svyazannom-s-trudovoy-deyatelnostyu/?
sphrase\_id=1588311. - Қaralғan kүnі: kүnі: 28.12.2025. {[}in Kazakh{]}
\end{refs}

\begin{info}
\hspace{1em}\emph{{\bfseries Авторлар туралы мәліметтер}}

Актаева Л.М.- медицина ғылымдарының докторы, Еңбекті қорғау жөніндегі
республикалық ғылыми-зерттеу институтының Бас директоры, Астана,
Қазақстан e-mail:rniiot@rniiot.kz;

Бекмағамбетов Ә.Б.- заң ғылымдарының кандидаты, қауымдастырылған
профессорі (доцент), Еңбекті қорғау жөніндегі республикалық
ғылыми-зерттеу институтының Бас директорының орынбасары, Кравцова көшесі
18, Астана e-mail:adilet1979@mail.ru;

Рахметова А.М. - медицина ғылымдарының кандидаты, доцент, ғылыми
зерттеулер орталығының бөлім басшысының орынбасары, Астана Қазақстан,
e-mail: ra\_anar@mail.ru;

Насырова Г.А.- экономика ғылымдарының докторы, ғылыми зерттеулер
орталығының жетекші ғылыми қызметкері, Астана Қазақстан, e-mail:
gnassyrova@yandex.kz;

Құлмағамбетова Э.А.- химия ғылымдарының кандидаты, ғылыми зерттеулер
орталығының жетекші ғылыми қызметкері, Астана, Қазақстан, e-mail:
elya\_kulmagambet@mail.ru;

Алшынбекова Г. Қ.- биология ғылымдарының кандидаты, ғылыми зерттеулер
орталығының жетекші ғылыми қызметкері, Астана Қазақстан, e-mail:
gulnaz\_gak@mail.ru;

Шилменова А.Т. - экономика ғылымдарының магистрі, ғылыми зерттеулер
орталығының жетекші ғылыми қызметкері, , Астана, Қазақстан, e-mail:
aigul.shilmenova@gmail.com;

Сағындықова Н. Т.- техника және технология ғылымдарының магистрі, ғылыми
зерттеулер орталығының аға ғылыми қызметкері, Астана, Қазақстан, e-mail:
nursag79@mail.ru.

\hspace{1em}\emph{{\bfseries Сведения об авторах}}

Актаева Л.М.- доктор медицинских наук, Генеральный директор
Республиканского научно-исследовательского института по охране труда, г.
Астана, Казахстан, e-mail: e-mail:rniiot@rniiot.kz;

Бекмагамбетов А.Б.-кандидат юридичесских наук, ассоциированный профессор
(доцент) заместитель генерального директора Республиканского
научно-исследовательского института по охране труда, Астана, Казахстан,
e-mail: e-mail:adilet1979@mail.ru;

Рахметова А.М.- кандидат медицинских наук, доцент, заместитель
руководителя отдела центра научных исследований, г. Астана, Казахстан,
e-mail: ra\_anar@mail.ru;

Насырова Г.А. - доктор экономических наук, руководитель отдела
стратегического развития и международного сотрудничества, Астана,
Казахстан, e-mail: gnassyrova@yandex.kz;

Кульмагамбетова Э.А.- кандидат химических наук, ведущий научный
сотрудник центра научных исследований, Астана, Казахстан, e-mail:
elya\_kulmagambet@mail.ru;

Алшынбекова Г.К-кандидат биологических наук, ведущий научный сотрудник
центра научных исследований, г. Астана, Казахстан, e-mail:
gulnaz\_gak@mail.ru;

Шилменова А.Т.- магистр экономических наук, старший научный сотрудник
центра научных исследований, ул. Кравцова 18, г. Астана, Казахстан,
e-mail:
aigul.shilmenova@gmail.com;

Сагиндикова Н.Т.-магистр техники и технологии, старший научный сотрудник
центра научных исследований, ул. Кравцова 18, г. Астана, Казахстан,
010000, e-mail: nursag79@mail.ru.
\end{info}
